\documentclass[11pt]{article}
\usepackage[a4paper,margin=25mm]{geometry}
\usepackage{amsmath,amssymb,amsthm}
\usepackage[hidelinks]{hyperref}
\emergencystretch=2em
\newtheorem{theorem}{Theorem}
\title{Cooperation discounts are not ratios of raw stage payoffs\\
\large Disproof of Conjecture 00000008240}
\author{GitHub: gaochengzhecpu}\date{}
\begin{document}\maketitle
\begin{abstract}
The first assertion of the conjecture identifies the critical discount
for cooperation with the ratio of the one-period defection payoff to
the cooperative payoff. A standard repeated prisoner's dilemma has a
cooperative subgame-perfect equilibrium at discount $1/2$, whereas that
ratio is $3/2$. We verify the equilibrium against every unilateral
infinite action stream after every finite history. Lean formalizes the
stage game, arbitrary pure history-dependent profiles, their infinite
discounted payoffs, subgame perfection, and the resulting inequality for
the infimum over all cooperative pure equilibria. The example and its
equilibrium argument are classical; no new folk theorem is claimed.
\end{abstract}

\section*{The assertion and its interpretation}
Both language versions assert the raw-payoff ratio
\[
 \delta_c=\frac{\text{one-period defection payoff}}
 {\text{cooperative payoff}}.
\]
We interpret cooperation as perpetual mutual cooperation on the
equilibrium path, with subgame-perfect incentives under perfect public
monitoring. For a given stage game let $\mathcal S$ be the set of
$\delta\in(0,1)$ for which such a pure history-dependent equilibrium
exists, and define $\delta_c=\inf\mathcal S$. We will exhibit a member
of $\mathcal S$ and prove that the set is bounded below, so this infimum
is well defined. Any definition allowing mixed equilibria as well has
a weakly smaller critical discount and gives the same contradiction.

The source does not specify a noise model or a particular family of
static profiles. We do not assign definitions to those later conjuncts:
failure of the first conjunct refutes the conjunction. In particular,
we do not interpret ``defection payoff'' as an already normalized
payoff \emph{difference}. The standard difference formula
$(T-R)/(T-P)$ for grim-trigger incentives is a different expression.

\section*{A standard stage game}
Consider two players with actions $C,D$ and payoffs
\[
\begin{array}{c|cc}
 & C & D\\\hline
 C &(2,2)&(0,3)\\
 D &(3,0)&(1,1)
\end{array}
\]
Thus $T=3>R=2>P=1>S=0$. Each player's infinite payoff is
$\sum_{n\geq0}\delta^n u_n$. Multiplying this by $1-\delta$ would
leave all incentives unchanged. Every such series converges for
$0\leq\delta<1$ because $0\leq u_n\leq3$.

This is a familiar textbook example, including the exact stage-payoff
table in the course handout \cite{osborne}. The cooperative
grim-trigger equilibrium at $\delta=1/2$ is also Proposition~1 of
Levin's notes \cite{levin}, after adding one to every stage payoff.
That addition preserves every discounted payoff comparison.

\begin{theorem}
This game has a subgame-perfect equilibrium with perpetual mutual
cooperation at $\delta=1/2$. Consequently $\delta_c\leq1/2$, whereas
$T/R=3/2$. The first assertion of Conjecture 00000008240 is false.
\end{theorem}
\begin{proof}
Both players use grim trigger: cooperate until a defection by either
player has occurred, and defect forever after that. At any finite
history the public state is either $c$, with no past defection, or
$p$, with punishment already triggered. At $\delta=1/2$, following
the strategy gives continuation values
\[
 V(c)=2\sum_{n\geq0}2^{-n}=4,
 \qquad V(p)=\sum_{n\geq0}2^{-n}=2.
\]
For an arbitrary deviator, the stage reward $r$, next state, and
one-step value satisfy the following exhaustive table:
\[
\begin{array}{c|c|c|c|c}
 \text{state}&\text{own action}&r&\text{next state}
     &r+\tfrac12 V(\text{next state})\\\hline
 c&C&2&c&4\\
 c&D&3&p&4\\
 p&C&0&p&1\\
 p&D&1&p&2
\end{array}
\]
In each case the last entry is at most the current state's value.
For any infinite unilateral action stream, induction and telescoping
therefore give, for every $N$,
\[
 \sum_{n=0}^{N-1}2^{-n}r_n+2^{-N}V(s_N)\leq V(s_0).
\]
The terminal value is nonnegative. Letting $N$ tend to infinity shows
that the deviator's total payoff is at most $V(s_0)$, exactly the
payoff from following grim trigger. This argument applies to both
players and every finite history, including histories off the
cooperative path. It proves subgame perfection against arbitrary
pure deviations, not just a single defection. The same pointwise
bound also rules out gains from randomized deviations by taking
expectations.

Starting from the empty history, both players cooperate forever.
Thus $1/2\in\mathcal S$ and $\mathcal S\subset(0,1)$, whence
$\delta_c\leq1/2$. Since the raw defection-to-cooperation payoff
ratio is $3/2>1/2$, the asserted equality fails.
\end{proof}

\section*{Why raw payoff ratios are structurally unsuitable}
Adding the same constant to every outcome payoff of a player adds a
strategy-independent constant $k/(1-\delta)$ to the repeated payoff.
It cannot change any equilibrium or the cooperation threshold. In
contrast, the raw ratio $T/R$ changes to $(T+k)/(R+k)$.
This is the basic defect in the proposed ratio. The counterexample
uses an ordinary strictly ordered prisoner's dilemma and a standard
cooperative equilibrium, with no degenerate discount or payoff.

\section*{Formal verification and scope}
The Lean namespace is \texttt{Conjecture8240}. Booleans encode the
two actions, and the function \texttt{stage} is the first player's
actual table; the second player's payoff is the same table with
arguments interchanged. Finite histories store both players' actions.
A \texttt{Profile} assigns an action to each player at every history.
The recursive trajectory against a fixed opponent generates the
actual history for each unilateral infinite action stream.

The definition \texttt{SPE} quantifies over both players, every finite
history, a compliant continuation, and every possible unilateral
infinite stream. The definition \texttt{Cooperates} requires the
all-$C$ stream to be compliant from the empty history for each player.
\texttt{sustainable} ranges over all such pure profiles at
$\delta\in(0,1)$; its infimum is \texttt{critical}. Thus the formal
critical value is not restricted to grim-trigger profiles.

The central formally verified results are
\[
\begin{gathered}
 \texttt{history\_payoff\_summable},\quad
 \texttt{grim\_is\_spe},\quad
 \texttt{grim\_cooperates},\\
 \texttt{half\_sustainable},\quad
 \texttt{critical\_le\_half},\quad
 \texttt{conjecture8240\_false}.
\end{gathered}
\]
The final theorem negates \texttt{ClaimedRatio}, the specialization
of the first conjectured equality to this explicit game.
All objects needed for the pure-strategy counterexample, including
convergence of discounted sums, are formalized. The paper's comments
about randomized deviations, the larger mixed-strategy infimum, and
invariance under additive payoff shifts are not separately formalized.
No exact optimal value for $\delta_c$ is needed or claimed here.
Only the usual Lean axioms of propositional extensionality, classical
choice, and quotient soundness occur in the printed axiom audit.

\section*{Source and reproduction}
The unmodified bilingual conjecture is included as \texttt{SOURCE.md},
retrieved from upstream commit
\begin{center}\small
\texttt{b355e42906e6eaf52bc7e9d005e0ef02d8aa7400}.
\end{center}
Its SHA-256 digest and source path appear in
\texttt{verification/SOURCE\_PROVENANCE.md}.
The project pins Lean 4.19.0 and Mathlib commit
\begin{center}\small
\texttt{c44e0c8ee63ca166450922a373c7409c5d26b00b}.
\end{center}
From the \texttt{lean} directory, run \texttt{lake build}, or
\begin{center}\small
\texttt{lake env lean -DwarningAsError=true Main.lean}.
\end{center}
The work was developed by a delegated AI agent with a self-review;
the coordinating agent reviews it separately. There is no external
independent review.

\begin{thebibliography}{9}
\bibitem{levin} Jonathan Levin, \emph{Repeated Games I: Perfect
Monitoring}, May 2006, Proposition~1.
\url{https://web.stanford.edu/~jdlevin/Econ%20286/Repeated%20Games%20I.pdf}
\bibitem{osborne} \emph{Chapter 14: Repeated Prisoner's Dilemma},
Northwestern University course handout discussing Osborne's text,
stage game and grim-trigger strategies, pp.~1--2.
\url{https://sites.math.northwestern.edu/~clark/364/handouts/osborne14.pdf}
\end{thebibliography}
\end{document}
